% Options for packages loaded elsewhere
% Options for packages loaded elsewhere
\PassOptionsToPackage{unicode}{hyperref}
\PassOptionsToPackage{hyphens}{url}
\PassOptionsToPackage{dvipsnames,svgnames,x11names}{xcolor}
%
\documentclass[
  letterpaper,
  DIV=11,
  numbers=noendperiod]{scrartcl}
\usepackage{xcolor}
\usepackage{amsmath,amssymb}
\setcounter{secnumdepth}{-\maxdimen} % remove section numbering
\usepackage{iftex}
\ifPDFTeX
  \usepackage[T1]{fontenc}
  \usepackage[utf8]{inputenc}
  \usepackage{textcomp} % provide euro and other symbols
\else % if luatex or xetex
  \usepackage{unicode-math} % this also loads fontspec
  \defaultfontfeatures{Scale=MatchLowercase}
  \defaultfontfeatures[\rmfamily]{Ligatures=TeX,Scale=1}
\fi
\usepackage{lmodern}
\ifPDFTeX\else
  % xetex/luatex font selection
\fi
% Use upquote if available, for straight quotes in verbatim environments
\IfFileExists{upquote.sty}{\usepackage{upquote}}{}
\IfFileExists{microtype.sty}{% use microtype if available
  \usepackage[]{microtype}
  \UseMicrotypeSet[protrusion]{basicmath} % disable protrusion for tt fonts
}{}
\makeatletter
\@ifundefined{KOMAClassName}{% if non-KOMA class
  \IfFileExists{parskip.sty}{%
    \usepackage{parskip}
  }{% else
    \setlength{\parindent}{0pt}
    \setlength{\parskip}{6pt plus 2pt minus 1pt}}
}{% if KOMA class
  \KOMAoptions{parskip=half}}
\makeatother
% Make \paragraph and \subparagraph free-standing
\makeatletter
\ifx\paragraph\undefined\else
  \let\oldparagraph\paragraph
  \renewcommand{\paragraph}{
    \@ifstar
      \xxxParagraphStar
      \xxxParagraphNoStar
  }
  \newcommand{\xxxParagraphStar}[1]{\oldparagraph*{#1}\mbox{}}
  \newcommand{\xxxParagraphNoStar}[1]{\oldparagraph{#1}\mbox{}}
\fi
\ifx\subparagraph\undefined\else
  \let\oldsubparagraph\subparagraph
  \renewcommand{\subparagraph}{
    \@ifstar
      \xxxSubParagraphStar
      \xxxSubParagraphNoStar
  }
  \newcommand{\xxxSubParagraphStar}[1]{\oldsubparagraph*{#1}\mbox{}}
  \newcommand{\xxxSubParagraphNoStar}[1]{\oldsubparagraph{#1}\mbox{}}
\fi
\makeatother

\usepackage{color}
\usepackage{fancyvrb}
\newcommand{\VerbBar}{|}
\newcommand{\VERB}{\Verb[commandchars=\\\{\}]}
\DefineVerbatimEnvironment{Highlighting}{Verbatim}{commandchars=\\\{\}}
% Add ',fontsize=\small' for more characters per line
\usepackage{framed}
\definecolor{shadecolor}{RGB}{241,243,245}
\newenvironment{Shaded}{\begin{snugshade}}{\end{snugshade}}
\newcommand{\AlertTok}[1]{\textcolor[rgb]{0.68,0.00,0.00}{#1}}
\newcommand{\AnnotationTok}[1]{\textcolor[rgb]{0.37,0.37,0.37}{#1}}
\newcommand{\AttributeTok}[1]{\textcolor[rgb]{0.40,0.45,0.13}{#1}}
\newcommand{\BaseNTok}[1]{\textcolor[rgb]{0.68,0.00,0.00}{#1}}
\newcommand{\BuiltInTok}[1]{\textcolor[rgb]{0.00,0.23,0.31}{#1}}
\newcommand{\CharTok}[1]{\textcolor[rgb]{0.13,0.47,0.30}{#1}}
\newcommand{\CommentTok}[1]{\textcolor[rgb]{0.37,0.37,0.37}{#1}}
\newcommand{\CommentVarTok}[1]{\textcolor[rgb]{0.37,0.37,0.37}{\textit{#1}}}
\newcommand{\ConstantTok}[1]{\textcolor[rgb]{0.56,0.35,0.01}{#1}}
\newcommand{\ControlFlowTok}[1]{\textcolor[rgb]{0.00,0.23,0.31}{\textbf{#1}}}
\newcommand{\DataTypeTok}[1]{\textcolor[rgb]{0.68,0.00,0.00}{#1}}
\newcommand{\DecValTok}[1]{\textcolor[rgb]{0.68,0.00,0.00}{#1}}
\newcommand{\DocumentationTok}[1]{\textcolor[rgb]{0.37,0.37,0.37}{\textit{#1}}}
\newcommand{\ErrorTok}[1]{\textcolor[rgb]{0.68,0.00,0.00}{#1}}
\newcommand{\ExtensionTok}[1]{\textcolor[rgb]{0.00,0.23,0.31}{#1}}
\newcommand{\FloatTok}[1]{\textcolor[rgb]{0.68,0.00,0.00}{#1}}
\newcommand{\FunctionTok}[1]{\textcolor[rgb]{0.28,0.35,0.67}{#1}}
\newcommand{\ImportTok}[1]{\textcolor[rgb]{0.00,0.46,0.62}{#1}}
\newcommand{\InformationTok}[1]{\textcolor[rgb]{0.37,0.37,0.37}{#1}}
\newcommand{\KeywordTok}[1]{\textcolor[rgb]{0.00,0.23,0.31}{\textbf{#1}}}
\newcommand{\NormalTok}[1]{\textcolor[rgb]{0.00,0.23,0.31}{#1}}
\newcommand{\OperatorTok}[1]{\textcolor[rgb]{0.37,0.37,0.37}{#1}}
\newcommand{\OtherTok}[1]{\textcolor[rgb]{0.00,0.23,0.31}{#1}}
\newcommand{\PreprocessorTok}[1]{\textcolor[rgb]{0.68,0.00,0.00}{#1}}
\newcommand{\RegionMarkerTok}[1]{\textcolor[rgb]{0.00,0.23,0.31}{#1}}
\newcommand{\SpecialCharTok}[1]{\textcolor[rgb]{0.37,0.37,0.37}{#1}}
\newcommand{\SpecialStringTok}[1]{\textcolor[rgb]{0.13,0.47,0.30}{#1}}
\newcommand{\StringTok}[1]{\textcolor[rgb]{0.13,0.47,0.30}{#1}}
\newcommand{\VariableTok}[1]{\textcolor[rgb]{0.07,0.07,0.07}{#1}}
\newcommand{\VerbatimStringTok}[1]{\textcolor[rgb]{0.13,0.47,0.30}{#1}}
\newcommand{\WarningTok}[1]{\textcolor[rgb]{0.37,0.37,0.37}{\textit{#1}}}

\usepackage{longtable,booktabs,array}
\newcounter{none} % for unnumbered tables
\usepackage{calc} % for calculating minipage widths
% Correct order of tables after \paragraph or \subparagraph
\usepackage{etoolbox}
\makeatletter
\patchcmd\longtable{\par}{\if@noskipsec\mbox{}\fi\par}{}{}
\makeatother
% Allow footnotes in longtable head/foot
\IfFileExists{footnotehyper.sty}{\usepackage{footnotehyper}}{\usepackage{footnote}}
\makesavenoteenv{longtable}
\usepackage{graphicx}
\makeatletter
\newsavebox\pandoc@box
\newcommand*\pandocbounded[1]{% scales image to fit in text height/width
  \sbox\pandoc@box{#1}%
  \Gscale@div\@tempa{\textheight}{\dimexpr\ht\pandoc@box+\dp\pandoc@box\relax}%
  \Gscale@div\@tempb{\linewidth}{\wd\pandoc@box}%
  \ifdim\@tempb\p@<\@tempa\p@\let\@tempa\@tempb\fi% select the smaller of both
  \ifdim\@tempa\p@<\p@\scalebox{\@tempa}{\usebox\pandoc@box}%
  \else\usebox{\pandoc@box}%
  \fi%
}
% Set default figure placement to htbp
\def\fps@figure{htbp}
\makeatother





\setlength{\emergencystretch}{3em} % prevent overfull lines

\providecommand{\tightlist}{%
  \setlength{\itemsep}{0pt}\setlength{\parskip}{0pt}}



 


\KOMAoption{captions}{tableheading}
\makeatletter
\@ifpackageloaded{caption}{}{\usepackage{caption}}
\AtBeginDocument{%
\ifdefined\contentsname
  \renewcommand*\contentsname{Table of contents}
\else
  \newcommand\contentsname{Table of contents}
\fi
\ifdefined\listfigurename
  \renewcommand*\listfigurename{List of Figures}
\else
  \newcommand\listfigurename{List of Figures}
\fi
\ifdefined\listtablename
  \renewcommand*\listtablename{List of Tables}
\else
  \newcommand\listtablename{List of Tables}
\fi
\ifdefined\figurename
  \renewcommand*\figurename{Figure}
\else
  \newcommand\figurename{Figure}
\fi
\ifdefined\tablename
  \renewcommand*\tablename{Table}
\else
  \newcommand\tablename{Table}
\fi
}
\@ifpackageloaded{float}{}{\usepackage{float}}
\floatstyle{ruled}
\@ifundefined{c@chapter}{\newfloat{codelisting}{h}{lop}}{\newfloat{codelisting}{h}{lop}[chapter]}
\floatname{codelisting}{Listing}
\newcommand*\listoflistings{\listof{codelisting}{List of Listings}}
\makeatother
\makeatletter
\makeatother
\makeatletter
\@ifpackageloaded{caption}{}{\usepackage{caption}}
\@ifpackageloaded{subcaption}{}{\usepackage{subcaption}}
\makeatother
\usepackage{bookmark}
\IfFileExists{xurl.sty}{\usepackage{xurl}}{} % add URL line breaks if available
\urlstyle{same}
\hypersetup{
  pdftitle={Week\_1},
  colorlinks=true,
  linkcolor={blue},
  filecolor={Maroon},
  citecolor={Blue},
  urlcolor={Blue},
  pdfcreator={LaTeX via pandoc}}


\title{Week\_1}
\author{}
\date{}
\begin{document}
\maketitle


\section{Introduction}\label{introduction}

Welcome! This first session covers the basics you need to learning R:
installing R and RStudio, organizing a project, understanding the
RStudio interface, and the core building blocks of the R language
(objects, functions, operators, and loops).

\subsection{Install R and R studio}\label{install-r-and-r-studio}

\begin{itemize}
\tightlist
\item
  \textbf{R} is the programming language and engine that actually runs
  your code.
\item
  \textbf{RStudio} is the interface (IDE) that makes working with R
  easier --- script editor, console, plots, help.
\end{itemize}

\begin{enumerate}
\def\labelenumi{\arabic{enumi}.}
\item
  Install R:
  \url{https://cran.r-project.org/bin/windows/base/R-4.6.1-win.exe}
\item
  Install RStudio: \url{https://posit.co/downloads}
\item
  R bible: \url{https://r4ds.hadley.nz/}, \emph{R for Data Science}
\end{enumerate}

\subsection{Project and folder organization --- good
practices}\label{project-and-folder-organization-good-practices}

\begin{enumerate}
\def\labelenumi{\arabic{enumi}.}
\tightlist
\item
  Create one general folder per research project
\item
  Split the folder into subfolders. A common structure (there's no
  single ``correct'' recipe, adapt to your workflow):

  \begin{itemize}
  \tightlist
  \item
    \texttt{data/} --- raw and processed data
  \item
    \texttt{scripts/} (or \texttt{R/}) --- your \texttt{.R} /
    \texttt{.qmd} scripts
  \item
    \texttt{outputs/} --- figures, tables, model results
  \item
    \texttt{papers/} or \texttt{docs/} --- literature, manuscripts,
    notes
  \end{itemize}
\item
  Good coding practices:

  \begin{itemize}
  \tightlist
  \item
    \textbf{Comment your code}. Future-you (and your collaborators) will
    thank you.
  \item
    \textbf{Use an outline/layout} --- section headers (\texttt{\#},
    \texttt{\#\#}, \texttt{\#\#\#} in comments or Markdown) keep long
    scripts navigable.
  \item
    \textbf{Use a consistent naming/writing standard} --- e.g.~always
    \texttt{snake\_case} for object names, consistent indentation,
    consistent spacing around \texttt{\textless{}-}.
  \end{itemize}
\end{enumerate}

\subsection{Interface R}\label{interface-r}

RStudio is organized into four main panes:

\begin{enumerate}
\def\labelenumi{\arabic{enumi}.}
\tightlist
\item
  \textbf{Source / R Script} (top-left) --- where you write and save
  code.
\item
  \textbf{Console} (bottom-left) --- where code actually runs and have
  the error messages.
\item
  \textbf{Environment / History} (top-right) --- shows the objects
  currently loaded in memory, and a record of past commands.
\item
  \textbf{Files / Plots / Packages / Help} (bottom-right) --- file
  browser, plot viewer, package manager, and documentation.
\end{enumerate}

\subsection{Logical programming}\label{logical-programming}

Programming means telling the computer a precise sequence of steps ---
think of it as a recipe. Before writing code, it helps to think through
the logic first: what are the inputs, what steps happen, what is the
output?

\subsection{Keyboard shortcuts}\label{keyboard-shortcuts}

{\def\LTcaptype{none} % do not increment counter
\begin{longtable}[]{@{}
  >{\raggedright\arraybackslash}p{(\linewidth - 2\tabcolsep) * \real{0.3662}}
  >{\raggedright\arraybackslash}p{(\linewidth - 2\tabcolsep) * \real{0.6338}}@{}}
\toprule\noalign{}
\begin{minipage}[b]{\linewidth}\raggedright
Shortcut
\end{minipage} & \begin{minipage}[b]{\linewidth}\raggedright
Action
\end{minipage} \\
\midrule\noalign{}
\endhead
\bottomrule\noalign{}
\endlastfoot
\textbf{Ctrl + Enter} & Run current line / selection \\
\textbf{Alt + -} & Insert \texttt{\textless{}-} \\
\textbf{Ctrl + Alt + I} & Insert a new code chunk (R Markdown/Quarto) \\
\textbf{Ctrl + Shift + M} & Insert pipe \texttt{\%\textgreater{}\%} \\
\textbf{Ctrl + Shift + C} & Comment/uncomment selected lines \\
\textbf{Ctrl + L} & Clear the console \\
\textbf{Ctrl + Shift + Enter} & Run entire chunk \\
\end{longtable}
}

\subsection{Operators}\label{operators}

\textbf{Arithmetic:} \texttt{+}, \texttt{-}, \texttt{/}, \texttt{*},
\texttt{\%*\%} (matrix multiplication), \texttt{\^{}} (power)

\textbf{Relational:} \texttt{\textgreater{}}, \texttt{\textless{}},
\texttt{\textgreater{}=}, \texttt{\textless{}=}, \texttt{==} (equal to),
\texttt{!=} (not equal to)

\textbf{Logical:} \texttt{\&} (and), \texttt{\textbar{}} (or),
\texttt{!} (not)

\begin{Shaded}
\begin{Highlighting}[]
\DecValTok{3} \SpecialCharTok{\^{}} \DecValTok{3}
\end{Highlighting}
\end{Shaded}

\begin{verbatim}
[1] 27
\end{verbatim}

\begin{Shaded}
\begin{Highlighting}[]
\DecValTok{2} \SpecialCharTok{*} \DecValTok{3}
\end{Highlighting}
\end{Shaded}

\begin{verbatim}
[1] 6
\end{verbatim}

\begin{Shaded}
\begin{Highlighting}[]
\DecValTok{2} \SpecialCharTok{+} \DecValTok{3}
\end{Highlighting}
\end{Shaded}

\begin{verbatim}
[1] 5
\end{verbatim}

\begin{Shaded}
\begin{Highlighting}[]
\DecValTok{2} \SpecialCharTok{{-}} \DecValTok{3}
\end{Highlighting}
\end{Shaded}

\begin{verbatim}
[1] -1
\end{verbatim}

\begin{Shaded}
\begin{Highlighting}[]
\DecValTok{2} \SpecialCharTok{/} \DecValTok{3}
\end{Highlighting}
\end{Shaded}

\begin{verbatim}
[1] 0.6666667
\end{verbatim}

\begin{Shaded}
\begin{Highlighting}[]
\DecValTok{3} \SpecialCharTok{\textgreater{}} \DecValTok{2}
\end{Highlighting}
\end{Shaded}

\begin{verbatim}
[1] TRUE
\end{verbatim}

\begin{Shaded}
\begin{Highlighting}[]
\DecValTok{3} \SpecialCharTok{==} \DecValTok{2}
\end{Highlighting}
\end{Shaded}

\begin{verbatim}
[1] FALSE
\end{verbatim}

\begin{Shaded}
\begin{Highlighting}[]
\ConstantTok{TRUE} \SpecialCharTok{\&} \ConstantTok{FALSE}   \CommentTok{\# AND {-}\textgreater{} both must be true}
\end{Highlighting}
\end{Shaded}

\begin{verbatim}
[1] FALSE
\end{verbatim}

\begin{Shaded}
\begin{Highlighting}[]
\ConstantTok{TRUE} \SpecialCharTok{|} \ConstantTok{FALSE}   \CommentTok{\# OR  {-}\textgreater{} at least one must be true}
\end{Highlighting}
\end{Shaded}

\begin{verbatim}
[1] TRUE
\end{verbatim}

\begin{Shaded}
\begin{Highlighting}[]
\SpecialCharTok{!}\ConstantTok{TRUE}          \CommentTok{\# NOT {-}\textgreater{} flips the value}
\end{Highlighting}
\end{Shaded}

\begin{verbatim}
[1] FALSE
\end{verbatim}

\section{Types of objects in R}\label{types-of-objects-in-r}

Good object names: short, clear, self-explanatory, no spaces, avoid
starting with numbers, avoid special characters (use \texttt{\_}
instead).

\begin{enumerate}
\def\labelenumi{\arabic{enumi}.}
\tightlist
\item
  \textbf{Integer} --- whole numbers
\item
  \textbf{Numeric} (double) --- decimal numbers
\item
  \textbf{Character} --- text/strings
\item
  \textbf{Logical} --- \texttt{TRUE} / \texttt{FALSE}
\item
  \textbf{Vector} --- a sequence of values\\
\item
  \textbf{Matrix} --- a 2-dimensional structure
\item
  \textbf{Data frame} --- a table with observation and variables
\item
  \textbf{List} --- a flexible container that can hold objects of
  different types and sizes together
\end{enumerate}

\begin{Shaded}
\begin{Highlighting}[]
\NormalTok{a }\OtherTok{\textless{}{-}} \DecValTok{1}
\NormalTok{a }\OtherTok{\textless{}{-}} \FunctionTok{as.integer}\NormalTok{(a)}
\FunctionTok{class}\NormalTok{(a)}
\end{Highlighting}
\end{Shaded}

\begin{verbatim}
[1] "integer"
\end{verbatim}

\begin{Shaded}
\begin{Highlighting}[]
\NormalTok{b }\OtherTok{\textless{}{-}} \FloatTok{2.3}
\FunctionTok{class}\NormalTok{(b)}
\end{Highlighting}
\end{Shaded}

\begin{verbatim}
[1] "numeric"
\end{verbatim}

\begin{Shaded}
\begin{Highlighting}[]
\NormalTok{c\_var }\OtherTok{\textless{}{-}} \FunctionTok{c}\NormalTok{(}\StringTok{"2"}\NormalTok{, }\StringTok{"3"}\NormalTok{)   }\CommentTok{\# character vector}
\FunctionTok{class}\NormalTok{(c\_var)}
\end{Highlighting}
\end{Shaded}

\begin{verbatim}
[1] "character"
\end{verbatim}

\begin{Shaded}
\begin{Highlighting}[]
\NormalTok{d }\OtherTok{\textless{}{-}} \ConstantTok{TRUE}
\FunctionTok{class}\NormalTok{(d)}
\end{Highlighting}
\end{Shaded}

\begin{verbatim}
[1] "logical"
\end{verbatim}

\begin{Shaded}
\begin{Highlighting}[]
\NormalTok{e }\OtherTok{\textless{}{-}} \FunctionTok{matrix}\NormalTok{(}\FunctionTok{c}\NormalTok{(}\DecValTok{1}\SpecialCharTok{:}\DecValTok{4}\NormalTok{), }\AttributeTok{ncol =} \DecValTok{2}\NormalTok{)}
\NormalTok{e}
\end{Highlighting}
\end{Shaded}

\begin{verbatim}
     [,1] [,2]
[1,]    1    3
[2,]    2    4
\end{verbatim}

\begin{Shaded}
\begin{Highlighting}[]
\NormalTok{f }\OtherTok{\textless{}{-}} \FunctionTok{data.frame}\NormalTok{(}
  \AttributeTok{yield    =} \FunctionTok{c}\NormalTok{(}\DecValTok{80}\NormalTok{, }\DecValTok{90}\NormalTok{, }\DecValTok{100}\NormalTok{),}
  \AttributeTok{N\_rate   =} \FunctionTok{c}\NormalTok{(}\DecValTok{20}\NormalTok{, }\DecValTok{50}\NormalTok{, }\DecValTok{80}\NormalTok{),}
  \AttributeTok{Location =} \StringTok{"Ashland"}\NormalTok{,}
  \AttributeTok{wheat    =} \FunctionTok{c}\NormalTok{(}\ConstantTok{TRUE}\NormalTok{, }\ConstantTok{TRUE}\NormalTok{, }\ConstantTok{FALSE}\NormalTok{)}
\NormalTok{)}
\NormalTok{f}
\end{Highlighting}
\end{Shaded}

\begin{verbatim}
  yield N_rate Location wheat
1    80     20  Ashland  TRUE
2    90     50  Ashland  TRUE
3   100     80  Ashland FALSE
\end{verbatim}

\begin{Shaded}
\begin{Highlighting}[]
\FunctionTok{str}\NormalTok{(f)}
\end{Highlighting}
\end{Shaded}

\begin{verbatim}
'data.frame':   3 obs. of  4 variables:
 $ yield   : num  80 90 100
 $ N_rate  : num  20 50 80
 $ Location: chr  "Ashland" "Ashland" "Ashland"
 $ wheat   : logi  TRUE TRUE FALSE
\end{verbatim}

\section{Function and packages}\label{function-and-packages}

\textbf{What is a function?} A set of instructions, packaged under a
name, that executes a sequence of operations on some input and returns a
result. Instead of retyping the same steps every time, you call the
function.

\textbf{Common built-in functions:}

\begin{enumerate}
\def\labelenumi{\arabic{enumi}.}
\tightlist
\item
  \texttt{sum()}
\item
  \texttt{read\_xlsx()} / \texttt{read.csv()}
\item
  \texttt{summary()}
\item
  \texttt{head()}, \texttt{tail()}
\item
  \texttt{str()}
\item
  \texttt{sqrt()}
\item
  \texttt{table()}
\item
  \texttt{seq()}
\item
  \texttt{length()}
\item
  \texttt{colnames()}
\item
  \texttt{print()}
\end{enumerate}

\begin{Shaded}
\begin{Highlighting}[]
\FunctionTok{c}\NormalTok{(}\DecValTok{1}\NormalTok{,}\DecValTok{2}\NormalTok{,}\DecValTok{3}\NormalTok{,}\DecValTok{4}\NormalTok{,}\DecValTok{5}\NormalTok{,}\DecValTok{6}\NormalTok{,}\DecValTok{7}\NormalTok{,}\DecValTok{8}\NormalTok{,}\DecValTok{9}\NormalTok{,}\DecValTok{10}\NormalTok{)}
\end{Highlighting}
\end{Shaded}

\begin{verbatim}
 [1]  1  2  3  4  5  6  7  8  9 10
\end{verbatim}

\begin{Shaded}
\begin{Highlighting}[]
\FunctionTok{c}\NormalTok{(}\DecValTok{1}\SpecialCharTok{:}\DecValTok{10}\NormalTok{)}
\end{Highlighting}
\end{Shaded}

\begin{verbatim}
 [1]  1  2  3  4  5  6  7  8  9 10
\end{verbatim}

\begin{Shaded}
\begin{Highlighting}[]
\FunctionTok{seq}\NormalTok{(}\DecValTok{1}\NormalTok{, }\DecValTok{10}\NormalTok{, }\DecValTok{1}\NormalTok{)}
\end{Highlighting}
\end{Shaded}

\begin{verbatim}
 [1]  1  2  3  4  5  6  7  8  9 10
\end{verbatim}

\textbf{What is a package?} A collection of functions (and sometimes
data) written by someone else, that you can install and load to expand
what R can do without needing anything else.

\textbf{Installing and using a package:}

\begin{Shaded}
\begin{Highlighting}[]
\CommentTok{\# install.packages("dplyr")   \# run once, to download the package}
\CommentTok{\# dplyr::filter()              \# use one function without loading the whole package}
\CommentTok{\# library(dplyr)               \# load the whole package for the session}
\end{Highlighting}
\end{Shaded}

\textbf{Main packages:}

\begin{enumerate}
\def\labelenumi{\arabic{enumi}.}
\tightlist
\item
  \textbf{base} --- pre-installed, core R functions
\item
  \textbf{ggplot2} --- data visualization
\item
  \textbf{dplyr} --- data manipulation/analysis
\item
  \textbf{readxl} --- reading Excel files
\item
  \textbf{tidyr} --- reshaping data (long/wide formats)
\end{enumerate}

\subsection{Writing your own function}\label{writing-your-own-function}

General structure:

\begin{Shaded}
\begin{Highlighting}[]
\NormalTok{new\_function }\OtherTok{\textless{}{-}} \ControlFlowTok{function}\NormalTok{(a, b, c) \{}

  \CommentTok{\# commands / calculations}

  \FunctionTok{return}\NormalTok{(result)}
\NormalTok{\}}
\end{Highlighting}
\end{Shaded}

Example:

\begin{Shaded}
\begin{Highlighting}[]
\NormalTok{my\_mean }\OtherTok{\textless{}{-}} \ControlFlowTok{function}\NormalTok{(a, b) \{}
\NormalTok{  result }\OtherTok{\textless{}{-}}\NormalTok{ (a }\SpecialCharTok{+}\NormalTok{ b) }\SpecialCharTok{/} \DecValTok{2}
  \FunctionTok{return}\NormalTok{(result)}
\NormalTok{\}}

\FunctionTok{my\_mean}\NormalTok{(}\DecValTok{2}\NormalTok{, }\DecValTok{4}\NormalTok{)}
\end{Highlighting}
\end{Shaded}

\begin{verbatim}
[1] 3
\end{verbatim}

\subsection{Loops}\label{loops}

Loops repeat a task a set number of times, or until a condition is no
longer true.

\begin{itemize}
\tightlist
\item
  \texttt{for()} --- repeat a fixed, known number of times
\item
  \texttt{while()} --- repeat as long as a condition stays true
\end{itemize}

Structure:

\begin{Shaded}
\begin{Highlighting}[]
\NormalTok{n }\OtherTok{=} \DecValTok{10}
\ControlFlowTok{for}\NormalTok{(i }\ControlFlowTok{in} \DecValTok{1}\SpecialCharTok{:}\NormalTok{n)\{}
  \FunctionTok{print}\NormalTok{(i)}
\NormalTok{\}}
\end{Highlighting}
\end{Shaded}

\begin{verbatim}
[1] 1
[1] 2
[1] 3
[1] 4
[1] 5
[1] 6
[1] 7
[1] 8
[1] 9
[1] 10
\end{verbatim}

\begin{Shaded}
\begin{Highlighting}[]
\NormalTok{i}\OtherTok{=} \DecValTok{1}
\ControlFlowTok{while}\NormalTok{(i }\SpecialCharTok{\textless{}}\NormalTok{ n) \{}
  
  \FunctionTok{print}\NormalTok{(i)}
  
\NormalTok{  i }\OtherTok{=}\NormalTok{ i }\SpecialCharTok{+} \DecValTok{1}
\NormalTok{\}}
\end{Highlighting}
\end{Shaded}

\begin{verbatim}
[1] 1
[1] 2
[1] 3
[1] 4
[1] 5
[1] 6
[1] 7
[1] 8
[1] 9
\end{verbatim}

\section{Conditions}\label{conditions}

\section{Practice exercises}\label{practice-exercises}

\begin{enumerate}
\def\labelenumi{\arabic{enumi}.}
\tightlist
\item
  Create a vector \texttt{n\_rates} with the values 0, 40, 80, 120, 160
  (kg N/ha).
\end{enumerate}

\begin{Shaded}
\begin{Highlighting}[]
\NormalTok{n\_rates }\OtherTok{=} \FunctionTok{c}\NormalTok{(}\DecValTok{0}\NormalTok{,}\DecValTok{40}\NormalTok{,}\DecValTok{80}\NormalTok{,}\DecValTok{120}\NormalTok{,}\DecValTok{160}\NormalTok{)}
\end{Highlighting}
\end{Shaded}

\begin{enumerate}
\def\labelenumi{\arabic{enumi}.}
\setcounter{enumi}{1}
\tightlist
\item
  Create a data frame with three columns: \texttt{plot} (1 to 5),
  \texttt{n\_rate} (\texttt{n\_rates} above), and \texttt{yield} (make
  up five plausible wheat yield values).
\end{enumerate}

\begin{Shaded}
\begin{Highlighting}[]
\FunctionTok{data.frame}\NormalTok{(}\AttributeTok{plot =} \FunctionTok{c}\NormalTok{(}\DecValTok{1}\SpecialCharTok{:}\DecValTok{5}\NormalTok{), }
           \AttributeTok{n\_rate =}\NormalTok{ n\_rates,}
           \AttributeTok{yield =} \FunctionTok{c}\NormalTok{(}\DecValTok{20}\NormalTok{,}\DecValTok{30}\NormalTok{,}\DecValTok{40}\NormalTok{,}\DecValTok{50}\NormalTok{,}\DecValTok{60}\NormalTok{))}
\end{Highlighting}
\end{Shaded}

\begin{verbatim}
  plot n_rate yield
1    1      0    20
2    2     40    30
3    3     80    40
4    4    120    50
5    5    160    60
\end{verbatim}

\begin{enumerate}
\def\labelenumi{\arabic{enumi}.}
\setcounter{enumi}{2}
\item
  Write a function \texttt{kg\_to\_lb(kg)} that converts kilograms to
  pounds (1 kg = 2.2046 lb) and test it on your \texttt{yield} column.
\item
  Write a \texttt{for} loop that prints ``Plot X received Y kg N/ha''
  for each row of your data frame.
\end{enumerate}




\end{document}
